% year: תשע"ט
% semester: א
% moed: א
% number: 4
% points: 20
% categories: func-analysis
% source: exams/מבחנים/תשעט/מועד א תשעט.pdf

\begin{question}
חקרו את הפונקציה $y=(x+2)e^{1/x}$ ושרטטו את הגרף שלה.

\underline{צריך לחקור}: תחום הגדרה, רציפות, אסימפטוטות, תחומי עליה וירידה, נקודות קיצון, תחומי קמירות וקעירות.
\end{question}

\begin{hint}
לאסימפטוטה המשופעת: $m=\lim\frac{y}{x}$, $n=\lim(y-mx)$; בחישוב $n$ הציבו $t=\frac1x$ והשתמשו ב-$\lim_{t\to0}\frac{e^t-1}{t}=1$.
\end{hint}

\begin{hint}
בדקו בנפרד את הגבולות החד-צדדיים ב-$x=0$: כאשר $x\to0^-$, $e^{1/x}\to0$.
\end{hint}

\begin{solution}
\textbf{תחום הגדרה ורציפות.} $x\neq0$. בתחום הפונקציה רציפה (מכפלה והרכבה של פונקציות רציפות). נקודת חיתוך עם ציר $x$: $x=-2$; $y>0$ עבור $x>-2$, $x\ne0$, ו-$y<0$ עבור $x<-2$.

\textbf{אסימפטוטות אנכיות.} כאשר $x\to0^+$: $\frac1x\to+\infty$, $e^{1/x}\to+\infty$ ו-$x+2\to2$, לכן $y\to+\infty$: הישר $x=0$ אסימפטוטה אנכית (מימין). כאשר $x\to0^-$: $e^{1/x}\to0$ ולכן $y\to0$ – משמאל אין אסימפטוטה (הגרף "נכנס" לראשית; ב-$x=0$ אי-רציפות עיקרית).

\textbf{אסימפטוטות משופעות.} $m=\lim_{x\to\pm\infty}\frac{(x+2)e^{1/x}}{x}=\lim\left(1+\frac2x\right)e^{1/x}=1$. כעת
\[ y-x=x\big(e^{1/x}-1\big)+2e^{1/x}=\frac{e^{1/x}-1}{1/x}+2e^{1/x}\xrightarrow[x\to\pm\infty]{}1+2=3, \]
לפי $\lim_{t\to0}\frac{e^t-1}{t}=1$ עם $t=\frac1x\to0$. לכן $y=x+3$ אסימפטוטה משופעת בשני הכיוונים.

\textbf{עליה וירידה.}
\[ y'=e^{1/x}+(x+2)e^{1/x}\cdot\left(-\frac1{x^2}\right)=e^{1/x}\,\frac{x^2-x-2}{x^2}=e^{1/x}\,\frac{(x-2)(x+1)}{x^2}. \]
סימן $y'$ כסימן $(x-2)(x+1)$: $y'>0$ ב-$(-\infty,-1)$ וב-$(2,\infty)$ – הפונקציה עולה; $y'<0$ ב-$(-1,0)$ וב-$(0,2)$ – הפונקציה יורדת.

\textbf{נקודות קיצון.} ב-$x=-1$ הנגזרת מחליפה סימן מ-$+$ ל-$-$: מקסימום מקומי $y(-1)=e^{-1}=\frac1e$. ב-$x=2$ מ-$-$ ל-$+$: מינימום מקומי $y(2)=4\sqrt e\approx6.59$.

\textbf{קמירות וקעירות.} חישוב ישיר (גזירת $e^{1/x}(1-\frac1x-\frac2{x^2})$) נותן
\[ y''=e^{1/x}\left(-\frac1{x^2}\right)\left(1-\frac1x-\frac2{x^2}\right)+e^{1/x}\left(\frac1{x^2}+\frac4{x^3}\right)=e^{1/x}\,\frac{5x+2}{x^4}. \]
סימן $y''$ כסימן $5x+2$: $y''<0$ (קעורה, $\cap$) ב-$(-\infty,-\frac25)$, ו-$y''>0$ (קמורה, $\cup$) ב-$(-\frac25,0)$ וב-$(0,\infty)$. נקודת פיתול: $x=-\frac25$, $y\left(-\frac25\right)=\frac85e^{-5/2}\approx0.131$.

\textbf{סקיצה.} מהפיתוח $y=x+3+\frac{5}{2x}+o\!\left(\frac1x\right)$ נובע שהגרף נמצא מתחת לאסימפטוטה כאשר $x\to-\infty$ ומעליה כאשר $x\to+\infty$. משמאל: הגרף מתקרב מלמטה לישר $y=x+3$, עולה וחוצה את ציר $x$ ב-$x=-2$, מגיע למקסימום $(-1,\frac1e)$, יורד דרך נקודת הפיתול $(-\frac25,\frac85e^{-5/2})$ ושואף ל-$0$ כאשר $x\to0^-$. מימין ל-$0$ הגרף יורד מ-$+\infty$ (אסימפטוטה $x=0$) עד המינימום $(2,4\sqrt e)$, ואז עולה ומתקרב לאסימפטוטה $y=x+3$ כאשר $x\to+\infty$.
\end{solution}
